\documentclass[10pt,a4paper]{amsart}
\usepackage[margin=19mm]{geometry}
\usepackage{amsmath,amssymb,mathtools}
\usepackage[scr=rsfs,cal=euler]{mathalfa}
\usepackage{xfrac}
\usepackage[unicode,hidelinks,bookmarks=false]{hyperref}
\newcommand{\ie}{i.e.\ }
\newcommand{\cf}{cf.\ }
\newcommand{\resp}{respectively\ }
\newcommand{\Subsection}{\subsection}
\providecommand{\nobreakdash}{\nobreak\hbox{-}}
\setlength{\emergencystretch}{2em}
\multlinegap0pt
\usepackage{bm}
\usepackage{bbm}
\usepackage{dsfont}
\usepackage{xspace}
\usepackage{cancel}
\usepackage{mathtools}
\usepackage{amsthm}
\makeatletter
\@ifundefined{theorem}{\newtheorem{theorem}{Theorem}}{}
\@ifundefined{lemma}{\newtheorem{lemma}[theorem]{Lemma}}{}
\makeatother
\def\cC{\mathcal{C}}
\def\dist{\mathrm{dist}}
\def\io{\int_{\Omega}}
\def\irn{\int_{\r^N}}
\def\n{\mathbb{N}}
\def\o{\Omega}
\def\rh{\rightharpoonup}
\def\rn{\mathbb{R}^N}
\def\r{\mathbb{R}}
\def\vp{\varphi}
\def\what{\widehat}
\title[Multiple nodal solutions to a scalar field equation with dou]{Multiple nodal solutions to a scalar field equation with double-power nonlinearity and zero mass at infinity}
\author{M\'onica Clapp, Carlos Culebro}
\date{Source: arXiv:2508.15167v1 (CC BY 4.0)}
\begin{document}
\maketitle

\noindent\emph{Source and licence.} Multiple nodal solutions to a scalar field equation with double-power nonlinearity and zero mass at infinity. Version arXiv:2508.15167v1 (CC BY 4.0), distributed under the Creative Commons Attribution 4.0 International licence, as recorded in the arXiv licence metadata for this exact version and shown on the abstract page https://arxiv.org/abs/2508.15167v1. Full rights notice: licenses/arxiv-2508.15167-CC-BY-4.0.txt.\par\smallskip

\noindent\textit{Editorial scope.} Counted unit: the Lemma whose source label is \texttt{lem:spliting} in the article (source0.tex lines 346--356, source label \texttt{lem:spliting}), delivered together with the article's own complete original proof of it (source0.tex lines 358--417, one \texttt{proof} environment of 60 lines). No mathematics is written, repaired or re-derived: statement and proof are the article's own text line for line. Required context retained uncounted: eq:limit problem (lines 196--199); lem:cm (lines 289--298); lem:lions (lines 318--322); lem:K (lines 330--338). Every definition, display and label of those lines is retained unchanged. Omitted from this package and not redistributed: 1--195, 200--288, 299--317, 323--329, 339--345, 357--357, 418--804. Nothing inside a retained excerpt is omitted or altered: every retained body is delivered exactly as the article's lines are written, so no omission receipt of any kind accompanies this package. Citations: the retained excerpts carry no citation command at all, so no bibliography entry of the article is transcribed and no reference is redistributed. Reference adaptation: none was needed and none was made; every reference inside the delivered unit resolves inside it, so no unresolved reference or citation is produced. Printed slips kept literal: no printed word, symbol, hypothesis or number of the counted statement or of its proof was altered. Layout and class: the article's own class is \texttt{article} with packages including inputenc, amsmath, amsfonts, amssymb, graphicx, verbatim, mathrsfs, upref, amsthm, amsxtra, exscale, cite, dsfont, hyperref; the wrapper is the lane's pinned \texttt{amsart} editorial wrapper at 10pt on A4, which restates the theorem-like environments the retained text needs and adds the article's own macro definitions that the retained text needs (\texttt{\textbackslash cC}, \texttt{\textbackslash dist}, \texttt{\textbackslash io}, \texttt{\textbackslash irn}, \texttt{\textbackslash n}, \texttt{\textbackslash o}, \texttt{\textbackslash r}, \texttt{\textbackslash rh}, \texttt{\textbackslash rn}, \texttt{\textbackslash vp}, \texttt{\textbackslash what}), with extra wrapper packages (bm, bbm, dsfont, xspace, cancel, mathtools). The delivered numbering is the unit's own. Assets: no retained span contains a figure environment, a graphics-inclusion command, a PDF-page inclusion, a table or a TikZ picture; no image file is copied into this package, the package declares no asset dependency and no image file is redistributed with it. Licence: version arXiv:2508.15167v1 of this article is distributed under the Creative Commons Attribution 4.0 International licence, as recorded in the arXiv licence metadata for this exact version and shown on the abstract page https://arxiv.org/abs/2508.15167v1. A full rights notice is delivered at licenses/arxiv-2508.15167-CC-BY-4.0.txt. Characters: every retained excerpt is pure ASCII, so the delivered engine, XeLaTeX with its default Latin Modern text font, reports no missing character.
\begin{equation}
\label{eq:limit problem}
-\Delta u = f(u), \qquad u\in D^{1,2}(\rn).
\end{equation}

\begin{lemma}\label{lem:cm} 
If $u_k\rh u$ weakly in $D^{1,2}(\rn)$ then, after passing to a subsequence, the following statements hold true:
\begin{itemize}
\item[$(a)$] $\|u_k\|_V^2 = \|u_k -u\|^2 + \|u\|_V^2 + o(1).$
\item[$(b)$] $\irn |f(u_k) - f(u)| |\vp|=o(1)$ \ for every $\vp\in\cC^{\infty}_{c}(\rn)$.
\item[$(c)$] $\irn F(u_k) = \irn F(u_k - u) + \irn F(u) + o(1)$.
\item[$(d)$] $\irn f(u_k)u_k = \irn f(u_k - u)(u_k-u) + \irn f(u)u + o(1)$.
\item[$(e)$] $Vu_k\to Vu$ \ and \ $f(u_k) - f(u_k -u) \to f(u)$  strongly in $(D^{1,2}_0(\o))'$.
\end{itemize}
\end{lemma}

\begin{lemma}\label{lem:lions} 
If $(u_k)$ is bounded in $D^{1,2}(\rn)$ and there exists $R>0$ such that 
$$\lim_{k\to\infty} \left( \sup_{y \in \rn} \int_{B_R(y)} |u_k|^2 \right)=0,$$
then $\lim_{k\to\infty} \irn f(u_k)u_k=0.$
\end{lemma}

\begin{lemma} \label{lem:K} 
Given a sequence $(y_k)$ in $\rn$ there exists a sequence $(\zeta_k)$ in $\rn$ and a closed subgroup $K$ of $G$ such that for some subsequence of $(y_k)$, denoted in the same way, the following hold:
\begin{itemize}
\item [$(a)$]The sequence $(\mathrm{dist} (Gy_k,\zeta_k))$ is bounded.
\item[$(b)$] $G_{\zeta_k}= K$ for all $k\in\n$.
\item[$(c)$] If $[G:K]<\infty$, then $|g\zeta_k - \what g \zeta_k|\to\infty$ for any $g,\what g\in G$ with $gK\neq \what g K$.
\item[$(d)$] If $[G:K]=\infty$ then, for any given $m\in\n$, there exist $g_1,\ldots,g_m\in G$ such that $|g_i\zeta_k - \what g_j\zeta_k|\to\infty$ if $i\neq j$.
\end{itemize}
\end{lemma}

\begin{lemma}\label{lem:spliting} 
Let $V=0$ and $(u_k)$ be a Palais-Smale sequence for $I_0$ in $D^{1,2}_0(\o)^G$ at the level $c$ such that  $u_k\rh 0$ weakly in $D^{1,2}_0(\o)$. 
\begin{itemize}
\item[$(a)$] If $\o$ is bounded, then $u_k\to 0$ strongly in $D^{1,2}_0(\o)$. 
\item[$(b)$] If $\o$ is an exterior domain and $(u_k)$ does not converge strongly to $0$ in $D^{1,2}_0(\o)$ then, after passing to a subsequence, there exist a sequence $(\zeta_k)$ in $\rn\smallsetminus\{0\}$, a closed subgroup $K$ of finite index in $G$, and a nontrivial solution $w$ to the limit problem \eqref{eq:limit problem} such that
$$G_{\zeta_k}= K\text{ \ for all \ }k \in \mathbb{N}\qquad\text{and}\qquad c\geq[G:K] \, I_{\infty}(w),$$ 
where $I_\infty:D^{1,2}(\rn)\to\r$,
$$I_\infty(w):=\frac{1}{2}\|w\|^2-\irn F(w),$$
is the energy functional of the limit problem \eqref{eq:limit problem}.  
\end{itemize}
\end{lemma}

\begin{proof}
Assume that $(u_k)$ does not converge strongly to $0$ in $D^{1,2}_0(\o)$. Then there exist $C_0>0$ and a subsequence such that $\|u_k\|\geq C_0$ for all $k\in\n$. As $(u_k)$ is bounded in $D^{1,2}_0(\o)$ and $I'_0(u_k)\to 0$, we have that
$$o(1)=I_0'(u_k)u_k= \|u_k\|^2 - \io f(u_k)u_k.$$
Hence,
$$0<C_0^2\leq\|u_k\|^2=\io f(u_k)u_k+o(1),$$
and, by Lemma \ref{lem:lions}, there exist $\delta >0$ and a sequence $(y_k)$ in $\rn$ such that
$$\int_{B_1(y_k)} |u_n|^2 = \sup_{y\in\rn} \int_{B_1(y)} |u_k|^2 \geq\delta.$$
For the sequence $(y_k)$, we fix a sequence $(\zeta_k)$ in $\rn$ and a closed subgroup $K$ of $G$ with the properties stated in Lemma \ref{lem:K}. In particular, there exist $g_k\in G$ and $C>0$ such that $|g_k^{-1}\zeta_k-y_k|=\mathrm{dist}(G\zeta_k,y_k)\leq C$. As $u_k$ is $G$-invariant, we get
\begin{equation}\label{eq:nontrivial}
\int_{B_{C+1}(\zeta_k)} |u_k|^{2} \geq \int_{B_1(g_k y_k)} |u_k|^{2} = \int_{B_1(y_k)} |u_k|^{2}\geq\delta>0.
\end{equation}
Set $w_k := u_k(\cdot + \zeta_k)$. Since $(w_k)$ is bounded in $D^{1,2}(\rn)$, passing to a subsequence, we have that
$w_k\rh w$ weakly in $D^{1,2}(\rn)$, $w_k\to w$ in $L^{2}_{\mathrm{loc}}(\rn)$ and $w_k\to w$ a.e. in $\rn$.
The inequality \eqref{eq:nontrivial} yields
$$\int_{B_{C+1}(0)} |w_k|^{2} \geq \delta>0.$$
Therefore, $w\neq 0$. Then, as $u_k \rh 0$ weakly in $D^{1,2}(\rn)$, an easy argument shows that $|\zeta_k|\to\infty$.

The inequality \eqref{eq:nontrivial} implies that $\dist(\zeta_k,\o)<C+1$ for all $k$. Thus, if $\o$ is bounded, then $(\zeta_k)$ is bounded and we obtain a contradiction. This proves $(a)$.

Assume now that $\o$ is an exterior domain. Given $\vp\in\cC^\infty_c(\rn)$ we set $\vp_k(x):=\vp(x-\zeta_k)$. Since $|\zeta_k|\to\infty$, we have that $\vp_k\in \cC^\infty_c(\o)$ for $k$ large enough. Then, using Lemma \ref{lem:cm}$(b)$ one sees that
$$I'_\infty(w)\vp=I'_\infty(w_k)\vp+o(1)=I'_0(u_k)\vp_k+o(1)=o(1).$$
Hence, $w$ is a nontrivial solution to the limit problem \eqref{eq:limit problem}.

Assume there exist $g_1,\ldots, g_m \in G$ such that $|g_i\zeta_k- g_j\zeta_k|\to\infty$ if $i \neq j$. Then, for each $j\in\{1,\ldots,m\}$,
$$g_j w_k - \sum_{i=j+1}^{m} (g_i w)(\cdot - g_i \zeta_k + g_j \zeta_k) \rh g_j  w \quad \text{weakly in \ } D^{1,2}(\rn),$$
with the sum being $0$ if $j=m$. Lemma \ref{lem:cm}$(c)$ yields
\begin{align*}
&\irn F\Big(g_jw_k - \sum_{i=j+1}^{m} (g_iw)(\cdot - g_i \zeta_k + g_j \zeta_k) \Big)\\
&\qquad= \irn F\Big(g_j w_k - \sum_{i=j}^{m} (g_iw)(\cdot - g_i \zeta_k + g_j \zeta_k)\Big) + \irn F(g_j w) + o(1).
\end{align*}
Since $u_k$  is $G$-invariant, performing the change of variable $x \mapsto x - g_j \zeta_k$ we derive
\begin{align*}
&\irn F\Big(u_k - \sum_{i=j+1}^{m} (g_i w)(\cdot - g_i \zeta_k ) \Big)= \irn F\Big(u_k - \sum_{i=j}^{m} (g_iw)(\cdot - g_i \zeta_k ) \Big) + \irn F(w) + o(1).
\end{align*}
Iterating this identity for $j= 1,\ldots, m$ we obtain
\begin{equation} \label{eq:F}
\irn F(u_k)=\irn F(u_k - z_k)+m\irn F(w) + o(1),
\end{equation}
where
$$z_k:=\sum_{i=1}^{m} (g_iw)(\cdot - g_i \zeta_k).$$
Similarly, using statements $(d)$ and $(a)$ of Lemma \ref{lem:cm} we obtain
\begin{equation} \label{eq:f}
\irn f(u_k)u_k=\irn f(u_k - z_k)(u_k - z_k)+m\irn f(w)w+o(1).
\end{equation}
and
\begin{equation}\label{eq:norm}
\|u_k\|^2 = \|u_k - z_k\|^2 + m \|w\|^2 + o(1).
\end{equation}
Since $w$ solves \eqref{eq:limit problem}, from \eqref{eq:norm} and \eqref{eq:f} we get
$$o(1)=I'_0(u_k)u_k=\|u_k-z_k\|^2-\irn f(u_k-z_k)(u_k-z_k)+o(1),$$
and from assumption $(f_2)$ we derive
$$\frac{1}{2}\|u_k-z_k\|^2-\irn F(u_k-z_k)=\irn\Big(\frac{1}{2}f(u_k-z_k)(u_k-z_k)-F(u_k-z_k)\Big)+o(1)\geq 0+o(1).$$
This inequality, combined with \eqref{eq:norm} and \eqref{eq:F}, yields
\begin{equation}\label{eq:lower bound}
c+o(1)=I_0(u_k)=\frac{1}{2}\|u_k-z_k\|^2-\irn F(u_k-z_k)+mI_\infty(w)+o(1)\geq mI_\infty(w)+o(1).
\end{equation}
Therefore, $m$ cannot be arbitrarily large. It follows from Lemma \ref{lem:K} that $[G:K]<\infty$. Setting $m:=[G:K]$ and passing to the limit in \eqref{eq:lower bound} we obtain
$$c\geq [G:K]\,I_\infty(w).$$
This completes the proof of $(b)$.
\end{proof}

\end{document}
